\documentclass[aspectratio=169]{beamer}
\usetheme{Madrid}
\usecolortheme{default}

\usepackage{amsmath,amssymb,amsfonts,bm,mathtools}
\usepackage{physics}

\title{Transformers: Mathematical Foundations}
\author{Morten Hjorth-Jensen}
\date{Spring 2026}

\begin{document}

\frame{\titlepage}

\begin{frame}{Structure of the lectures}
This lecture series is divided into four parts:

\begin{enumerate}
\item Foundations of sequence modeling and attention
\item Mathematical structure of self-attention and transformer blocks
\item Training, optimization, and scaling behavior
\item Theoretical interpretations and modern extensions
\end{enumerate}

The emphasis is on mathematical structure, linear algebra,
probability, and geometric interpretation of attention mechanisms.
\end{frame}

\section{Lecture 1: Sequence Modeling and Attention}


\begin{frame}{Probabilistic sequence modeling}

Given tokens $x_1,\dots,x_T$, the joint distribution factorizes via the chain rule

\[
p(x_1,\dots,x_T) = \prod_{t=1}^T p(x_t|x_1,\dots,x_{t-1}).
\]

Language models approximate the conditional distribution

\[
p_\theta(x_t|x_{<t}).
\]

\end{frame}


\begin{frame}{Token embeddings}

Each token $x$ in a vocabulary $\mathcal V$ is mapped to a vector

\[
e(x) \in \mathbb{R}^d.
\]

Embedding matrix

\[
E \in \mathbb{R}^{|\mathcal V|\times d}
\]

defines the representation space used by the model.

\end{frame}


\begin{frame}{Context representations}

Sequence models compute a contextual representation

\[
h_t = f(e_1,\dots,e_t)
\]

which summarizes the preceding tokens and is used to predict $x_{t+1}$.

\end{frame}


\begin{frame}{Limitations of recurrent models}

RNNs compute

\[
h_t = f(h_{t-1},x_t).
\]

Problems:

\begin{itemize}
\item sequential computation prevents parallelization
\item vanishing and exploding gradients
\item difficulty capturing long-range dependencies
\end{itemize}

\end{frame}


\begin{frame}{Attention as weighted aggregation}

Attention computes

\[
y = \sum_i \alpha_i v_i
\]

where $\alpha_i$ depend on similarity between a query $q$
and keys $k_i$.

\end{frame}


\begin{frame}{Softmax attention weights}

Scores

\[
s_i = q^T k_i
\]

Weights

\[
\alpha_i = \frac{\exp(s_i)}{\sum_j \exp(s_j)}
\]

ensure positive normalized contributions.

\end{frame}


\begin{frame}{Matrix form of attention}

For matrices

\[
Q,K,V,
\]

attention becomes

\[
\text{Attn}(Q,K,V) = \text{softmax}(QK^T)V.
\]

\end{frame}


\begin{frame}{Interpretation of attention}

Attention acts as

\begin{itemize}
\item similarity search
\item differentiable memory lookup
\item adaptive aggregation of contextual information
\end{itemize}

\end{frame}


\begin{frame}{Properties of attention (1)}

Consider the attention matrix

\[
A = \text{softmax}(QK^T).
\]

Properties:

\begin{itemize}
\item rows sum to one
\item defines a stochastic matrix
\item performs weighted mixing of value vectors
\end{itemize}

Output:

\[
Y = AV.
\]

\end{frame}


\begin{frame}{Properties of attention (2)}

Consider the attention matrix

\[
A = \text{softmax}(QK^T).
\]

Properties:

\begin{itemize}
\item rows sum to one
\item defines a stochastic matrix
\item performs weighted mixing of value vectors
\end{itemize}

Output:

\[
Y = AV.
\]

\end{frame}


\begin{frame}{Properties of attention (3)}

Consider the attention matrix

\[
A = \text{softmax}(QK^T).
\]

Properties:

\begin{itemize}
\item rows sum to one
\item defines a stochastic matrix
\item performs weighted mixing of value vectors
\end{itemize}

Output:

\[
Y = AV.
\]

\end{frame}


\begin{frame}{Properties of attention (4)}

Consider the attention matrix

\[
A = \text{softmax}(QK^T).
\]

Properties:

\begin{itemize}
\item rows sum to one
\item defines a stochastic matrix
\item performs weighted mixing of value vectors
\end{itemize}

Output:

\[
Y = AV.
\]

\end{frame}


\begin{frame}{Properties of attention (5)}

Consider the attention matrix

\[
A = \text{softmax}(QK^T).
\]

Properties:

\begin{itemize}
\item rows sum to one
\item defines a stochastic matrix
\item performs weighted mixing of value vectors
\end{itemize}

Output:

\[
Y = AV.
\]

\end{frame}


\begin{frame}{Properties of attention (6)}

Consider the attention matrix

\[
A = \text{softmax}(QK^T).
\]

Properties:

\begin{itemize}
\item rows sum to one
\item defines a stochastic matrix
\item performs weighted mixing of value vectors
\end{itemize}

Output:

\[
Y = AV.
\]

\end{frame}


\begin{frame}{Properties of attention (7)}

Consider the attention matrix

\[
A = \text{softmax}(QK^T).
\]

Properties:

\begin{itemize}
\item rows sum to one
\item defines a stochastic matrix
\item performs weighted mixing of value vectors
\end{itemize}

Output:

\[
Y = AV.
\]

\end{frame}


\begin{frame}{Properties of attention (8)}

Consider the attention matrix

\[
A = \text{softmax}(QK^T).
\]

Properties:

\begin{itemize}
\item rows sum to one
\item defines a stochastic matrix
\item performs weighted mixing of value vectors
\end{itemize}

Output:

\[
Y = AV.
\]

\end{frame}


\begin{frame}{Properties of attention (9)}

Consider the attention matrix

\[
A = \text{softmax}(QK^T).
\]

Properties:

\begin{itemize}
\item rows sum to one
\item defines a stochastic matrix
\item performs weighted mixing of value vectors
\end{itemize}

Output:

\[
Y = AV.
\]

\end{frame}


\begin{frame}{Properties of attention (10)}

Consider the attention matrix

\[
A = \text{softmax}(QK^T).
\]

Properties:

\begin{itemize}
\item rows sum to one
\item defines a stochastic matrix
\item performs weighted mixing of value vectors
\end{itemize}

Output:

\[
Y = AV.
\]

\end{frame}


\begin{frame}{Properties of attention (11)}

Consider the attention matrix

\[
A = \text{softmax}(QK^T).
\]

Properties:

\begin{itemize}
\item rows sum to one
\item defines a stochastic matrix
\item performs weighted mixing of value vectors
\end{itemize}

Output:

\[
Y = AV.
\]

\end{frame}


\begin{frame}{Properties of attention (12)}

Consider the attention matrix

\[
A = \text{softmax}(QK^T).
\]

Properties:

\begin{itemize}
\item rows sum to one
\item defines a stochastic matrix
\item performs weighted mixing of value vectors
\end{itemize}

Output:

\[
Y = AV.
\]

\end{frame}


\begin{frame}{Properties of attention (13)}

Consider the attention matrix

\[
A = \text{softmax}(QK^T).
\]

Properties:

\begin{itemize}
\item rows sum to one
\item defines a stochastic matrix
\item performs weighted mixing of value vectors
\end{itemize}

Output:

\[
Y = AV.
\]

\end{frame}


\begin{frame}{Properties of attention (14)}

Consider the attention matrix

\[
A = \text{softmax}(QK^T).
\]

Properties:

\begin{itemize}
\item rows sum to one
\item defines a stochastic matrix
\item performs weighted mixing of value vectors
\end{itemize}

Output:

\[
Y = AV.
\]

\end{frame}


\begin{frame}{Properties of attention (15)}

Consider the attention matrix

\[
A = \text{softmax}(QK^T).
\]

Properties:

\begin{itemize}
\item rows sum to one
\item defines a stochastic matrix
\item performs weighted mixing of value vectors
\end{itemize}

Output:

\[
Y = AV.
\]

\end{frame}


\section{Lecture 2: Self-Attention and Transformer Architecture}


\begin{frame}{Self-attention definition}

In self-attention the same sequence produces queries, keys and values:

\[
Q = XW_Q, \quad K = XW_K, \quad V = XW_V.
\]

Here

\[
X \in \mathbb{R}^{T\times d}.
\]

\end{frame}


\begin{frame}{Scaled dot-product attention}

Transformer attention uses

\[
\text{Attn}(Q,K,V)
=
\text{softmax}\left(\frac{QK^T}{\sqrt d}\right)V.
\]

Scaling prevents extremely large inner products.

\end{frame}


\begin{frame}{Dimensions}

If

\[
Q,K,V \in \mathbb{R}^{T\times d_k},
\]

then

\[
QK^T \in \mathbb{R}^{T\times T}.
\]

\end{frame}


\begin{frame}{Computational complexity}

Self-attention requires

\[
O(T^2 d)
\]

operations due to pairwise interactions between tokens.

\end{frame}


\begin{frame}{Multi-head attention}

Multiple attention heads

\[
\text{head}_i = \text{Attn}(Q_i,K_i,V_i).
\]

Combined output

\[
\text{Concat}(\text{head}_1,\dots,\text{head}_h)W_O.
\]

\end{frame}


\begin{frame}{Feed-forward network}

Each transformer layer includes

\[
\text{FFN}(x)=\sigma(xW_1+b_1)W_2+b_2.
\]

\end{frame}


\begin{frame}{Residual connections}

Residual updates

\[
x_{l+1} = x_l + f(x_l).
\]

This helps gradient flow in deep networks.

\end{frame}


\begin{frame}{Layer normalization}

Layer normalization rescales activations

\[
\text{LN}(x) = \frac{x-\mu}{\sigma}.
\]

This stabilizes training.

\end{frame}


\begin{frame}{Transformer block analysis (1)}

A transformer block contains:

\begin{enumerate}
\item multi-head self-attention
\item feed-forward network
\end{enumerate}

Both components operate on token representations and are
combined through residual connections.

\end{frame}


\begin{frame}{Transformer block analysis (2)}

A transformer block contains:

\begin{enumerate}
\item multi-head self-attention
\item feed-forward network
\end{enumerate}

Both components operate on token representations and are
combined through residual connections.

\end{frame}


\begin{frame}{Transformer block analysis (3)}

A transformer block contains:

\begin{enumerate}
\item multi-head self-attention
\item feed-forward network
\end{enumerate}

Both components operate on token representations and are
combined through residual connections.

\end{frame}


\begin{frame}{Transformer block analysis (4)}

A transformer block contains:

\begin{enumerate}
\item multi-head self-attention
\item feed-forward network
\end{enumerate}

Both components operate on token representations and are
combined through residual connections.

\end{frame}


\begin{frame}{Transformer block analysis (5)}

A transformer block contains:

\begin{enumerate}
\item multi-head self-attention
\item feed-forward network
\end{enumerate}

Both components operate on token representations and are
combined through residual connections.

\end{frame}


\begin{frame}{Transformer block analysis (6)}

A transformer block contains:

\begin{enumerate}
\item multi-head self-attention
\item feed-forward network
\end{enumerate}

Both components operate on token representations and are
combined through residual connections.

\end{frame}


\begin{frame}{Transformer block analysis (7)}

A transformer block contains:

\begin{enumerate}
\item multi-head self-attention
\item feed-forward network
\end{enumerate}

Both components operate on token representations and are
combined through residual connections.

\end{frame}


\begin{frame}{Transformer block analysis (8)}

A transformer block contains:

\begin{enumerate}
\item multi-head self-attention
\item feed-forward network
\end{enumerate}

Both components operate on token representations and are
combined through residual connections.

\end{frame}


\begin{frame}{Transformer block analysis (9)}

A transformer block contains:

\begin{enumerate}
\item multi-head self-attention
\item feed-forward network
\end{enumerate}

Both components operate on token representations and are
combined through residual connections.

\end{frame}


\begin{frame}{Transformer block analysis (10)}

A transformer block contains:

\begin{enumerate}
\item multi-head self-attention
\item feed-forward network
\end{enumerate}

Both components operate on token representations and are
combined through residual connections.

\end{frame}


\begin{frame}{Transformer block analysis (11)}

A transformer block contains:

\begin{enumerate}
\item multi-head self-attention
\item feed-forward network
\end{enumerate}

Both components operate on token representations and are
combined through residual connections.

\end{frame}


\begin{frame}{Transformer block analysis (12)}

A transformer block contains:

\begin{enumerate}
\item multi-head self-attention
\item feed-forward network
\end{enumerate}

Both components operate on token representations and are
combined through residual connections.

\end{frame}


\begin{frame}{Transformer block analysis (13)}

A transformer block contains:

\begin{enumerate}
\item multi-head self-attention
\item feed-forward network
\end{enumerate}

Both components operate on token representations and are
combined through residual connections.

\end{frame}


\begin{frame}{Transformer block analysis (14)}

A transformer block contains:

\begin{enumerate}
\item multi-head self-attention
\item feed-forward network
\end{enumerate}

Both components operate on token representations and are
combined through residual connections.

\end{frame}


\begin{frame}{Transformer block analysis (15)}

A transformer block contains:

\begin{enumerate}
\item multi-head self-attention
\item feed-forward network
\end{enumerate}

Both components operate on token representations and are
combined through residual connections.

\end{frame}


\begin{frame}{Transformer block analysis (16)}

A transformer block contains:

\begin{enumerate}
\item multi-head self-attention
\item feed-forward network
\end{enumerate}

Both components operate on token representations and are
combined through residual connections.

\end{frame}


\begin{frame}{Transformer block analysis (17)}

A transformer block contains:

\begin{enumerate}
\item multi-head self-attention
\item feed-forward network
\end{enumerate}

Both components operate on token representations and are
combined through residual connections.

\end{frame}


\begin{frame}{Transformer block analysis (18)}

A transformer block contains:

\begin{enumerate}
\item multi-head self-attention
\item feed-forward network
\end{enumerate}

Both components operate on token representations and are
combined through residual connections.

\end{frame}


\begin{frame}{Transformer block analysis (19)}

A transformer block contains:

\begin{enumerate}
\item multi-head self-attention
\item feed-forward network
\end{enumerate}

Both components operate on token representations and are
combined through residual connections.

\end{frame}


\begin{frame}{Transformer block analysis (20)}

A transformer block contains:

\begin{enumerate}
\item multi-head self-attention
\item feed-forward network
\end{enumerate}

Both components operate on token representations and are
combined through residual connections.

\end{frame}


\section{Lecture 3: Training Transformers}


\begin{frame}{Language modeling objective}

Training maximizes

\[
\sum_{t=1}^T \log p_\theta(x_t|x_{<t}).
\]

\end{frame}


\begin{frame}{Cross entropy loss}

Predicted distribution

\[
p_\theta = \text{softmax}(Wh).
\]

Loss

\[
\mathcal L = -\log p_\theta(x_t).
\]

\end{frame}


\begin{frame}{Gradient computation}

Gradients propagate through

\begin{itemize}
\item softmax normalization
\item matrix multiplications
\item linear projections
\end{itemize}

\end{frame}


\begin{frame}{Adam optimizer}

Transformer models commonly use

\begin{itemize}
\item Adam optimizer
\item learning rate warmup
\item weight decay
\end{itemize}

\end{frame}


\begin{frame}{Scaling laws}

Empirical scaling relationships

\[
\text{loss} \propto N^{-\alpha}.
\]

Performance improves with model size and dataset size.

\end{frame}


\begin{frame}{Optimization considerations (1)}

Training large transformers requires

\begin{itemize}
\item careful learning rate scheduling
\item gradient clipping
\item distributed training
\end{itemize}

Large models rely heavily on parallel hardware.

\end{frame}


\begin{frame}{Optimization considerations (2)}

Training large transformers requires

\begin{itemize}
\item careful learning rate scheduling
\item gradient clipping
\item distributed training
\end{itemize}

Large models rely heavily on parallel hardware.

\end{frame}


\begin{frame}{Optimization considerations (3)}

Training large transformers requires

\begin{itemize}
\item careful learning rate scheduling
\item gradient clipping
\item distributed training
\end{itemize}

Large models rely heavily on parallel hardware.

\end{frame}


\begin{frame}{Optimization considerations (4)}

Training large transformers requires

\begin{itemize}
\item careful learning rate scheduling
\item gradient clipping
\item distributed training
\end{itemize}

Large models rely heavily on parallel hardware.

\end{frame}


\begin{frame}{Optimization considerations (5)}

Training large transformers requires

\begin{itemize}
\item careful learning rate scheduling
\item gradient clipping
\item distributed training
\end{itemize}

Large models rely heavily on parallel hardware.

\end{frame}


\begin{frame}{Optimization considerations (6)}

Training large transformers requires

\begin{itemize}
\item careful learning rate scheduling
\item gradient clipping
\item distributed training
\end{itemize}

Large models rely heavily on parallel hardware.

\end{frame}


\begin{frame}{Optimization considerations (7)}

Training large transformers requires

\begin{itemize}
\item careful learning rate scheduling
\item gradient clipping
\item distributed training
\end{itemize}

Large models rely heavily on parallel hardware.

\end{frame}


\begin{frame}{Optimization considerations (8)}

Training large transformers requires

\begin{itemize}
\item careful learning rate scheduling
\item gradient clipping
\item distributed training
\end{itemize}

Large models rely heavily on parallel hardware.

\end{frame}


\begin{frame}{Optimization considerations (9)}

Training large transformers requires

\begin{itemize}
\item careful learning rate scheduling
\item gradient clipping
\item distributed training
\end{itemize}

Large models rely heavily on parallel hardware.

\end{frame}


\begin{frame}{Optimization considerations (10)}

Training large transformers requires

\begin{itemize}
\item careful learning rate scheduling
\item gradient clipping
\item distributed training
\end{itemize}

Large models rely heavily on parallel hardware.

\end{frame}


\begin{frame}{Optimization considerations (11)}

Training large transformers requires

\begin{itemize}
\item careful learning rate scheduling
\item gradient clipping
\item distributed training
\end{itemize}

Large models rely heavily on parallel hardware.

\end{frame}


\begin{frame}{Optimization considerations (12)}

Training large transformers requires

\begin{itemize}
\item careful learning rate scheduling
\item gradient clipping
\item distributed training
\end{itemize}

Large models rely heavily on parallel hardware.

\end{frame}


\begin{frame}{Optimization considerations (13)}

Training large transformers requires

\begin{itemize}
\item careful learning rate scheduling
\item gradient clipping
\item distributed training
\end{itemize}

Large models rely heavily on parallel hardware.

\end{frame}


\begin{frame}{Optimization considerations (14)}

Training large transformers requires

\begin{itemize}
\item careful learning rate scheduling
\item gradient clipping
\item distributed training
\end{itemize}

Large models rely heavily on parallel hardware.

\end{frame}


\begin{frame}{Optimization considerations (15)}

Training large transformers requires

\begin{itemize}
\item careful learning rate scheduling
\item gradient clipping
\item distributed training
\end{itemize}

Large models rely heavily on parallel hardware.

\end{frame}


\begin{frame}{Optimization considerations (16)}

Training large transformers requires

\begin{itemize}
\item careful learning rate scheduling
\item gradient clipping
\item distributed training
\end{itemize}

Large models rely heavily on parallel hardware.

\end{frame}


\begin{frame}{Optimization considerations (17)}

Training large transformers requires

\begin{itemize}
\item careful learning rate scheduling
\item gradient clipping
\item distributed training
\end{itemize}

Large models rely heavily on parallel hardware.

\end{frame}


\begin{frame}{Optimization considerations (18)}

Training large transformers requires

\begin{itemize}
\item careful learning rate scheduling
\item gradient clipping
\item distributed training
\end{itemize}

Large models rely heavily on parallel hardware.

\end{frame}


\begin{frame}{Optimization considerations (19)}

Training large transformers requires

\begin{itemize}
\item careful learning rate scheduling
\item gradient clipping
\item distributed training
\end{itemize}

Large models rely heavily on parallel hardware.

\end{frame}


\begin{frame}{Optimization considerations (20)}

Training large transformers requires

\begin{itemize}
\item careful learning rate scheduling
\item gradient clipping
\item distributed training
\end{itemize}

Large models rely heavily on parallel hardware.

\end{frame}


\section{Lecture 4: Theory and Extensions}


\begin{frame}{Attention as kernel regression}

Attention resembles kernel regression

\[
\alpha_i \propto \exp(q^T k_i).
\]

This acts as a similarity kernel between tokens.

\end{frame}


\begin{frame}{Attention matrix interpretation}

Attention defines a matrix

\[
A = \text{softmax}(QK^T).
\]

Then

\[
Y = AV.
\]

\end{frame}


\begin{frame}{Graph interpretation}

Tokens can be viewed as nodes in a graph
with edges weighted by attention coefficients.

\end{frame}


\begin{frame}{Transformers as dynamical systems}

Deep transformers can be written as

\[
x_{l+1} = x_l + f_l(x_l).
\]

This resembles discretized dynamical systems.

\end{frame}


\begin{frame}{Large scale transformers}

Modern transformer models contain billions of parameters
and are trained on trillions of tokens.

\end{frame}


\begin{frame}{Advanced transformer interpretation (1)}

Transformers can be analyzed through:

\begin{itemize}
\item kernel methods
\item graph neural network analogies
\item dynamical systems
\item spectral analysis of attention matrices
\end{itemize}

\end{frame}


\begin{frame}{Advanced transformer interpretation (2)}

Transformers can be analyzed through:

\begin{itemize}
\item kernel methods
\item graph neural network analogies
\item dynamical systems
\item spectral analysis of attention matrices
\end{itemize}

\end{frame}


\begin{frame}{Advanced transformer interpretation (3)}

Transformers can be analyzed through:

\begin{itemize}
\item kernel methods
\item graph neural network analogies
\item dynamical systems
\item spectral analysis of attention matrices
\end{itemize}

\end{frame}


\begin{frame}{Advanced transformer interpretation (4)}

Transformers can be analyzed through:

\begin{itemize}
\item kernel methods
\item graph neural network analogies
\item dynamical systems
\item spectral analysis of attention matrices
\end{itemize}

\end{frame}


\begin{frame}{Advanced transformer interpretation (5)}

Transformers can be analyzed through:

\begin{itemize}
\item kernel methods
\item graph neural network analogies
\item dynamical systems
\item spectral analysis of attention matrices
\end{itemize}

\end{frame}


\begin{frame}{Advanced transformer interpretation (6)}

Transformers can be analyzed through:

\begin{itemize}
\item kernel methods
\item graph neural network analogies
\item dynamical systems
\item spectral analysis of attention matrices
\end{itemize}

\end{frame}


\begin{frame}{Advanced transformer interpretation (7)}

Transformers can be analyzed through:

\begin{itemize}
\item kernel methods
\item graph neural network analogies
\item dynamical systems
\item spectral analysis of attention matrices
\end{itemize}

\end{frame}


\begin{frame}{Advanced transformer interpretation (8)}

Transformers can be analyzed through:

\begin{itemize}
\item kernel methods
\item graph neural network analogies
\item dynamical systems
\item spectral analysis of attention matrices
\end{itemize}

\end{frame}


\begin{frame}{Advanced transformer interpretation (9)}

Transformers can be analyzed through:

\begin{itemize}
\item kernel methods
\item graph neural network analogies
\item dynamical systems
\item spectral analysis of attention matrices
\end{itemize}

\end{frame}


\begin{frame}{Advanced transformer interpretation (10)}

Transformers can be analyzed through:

\begin{itemize}
\item kernel methods
\item graph neural network analogies
\item dynamical systems
\item spectral analysis of attention matrices
\end{itemize}

\end{frame}


\begin{frame}{Advanced transformer interpretation (11)}

Transformers can be analyzed through:

\begin{itemize}
\item kernel methods
\item graph neural network analogies
\item dynamical systems
\item spectral analysis of attention matrices
\end{itemize}

\end{frame}


\begin{frame}{Advanced transformer interpretation (12)}

Transformers can be analyzed through:

\begin{itemize}
\item kernel methods
\item graph neural network analogies
\item dynamical systems
\item spectral analysis of attention matrices
\end{itemize}

\end{frame}


\begin{frame}{Advanced transformer interpretation (13)}

Transformers can be analyzed through:

\begin{itemize}
\item kernel methods
\item graph neural network analogies
\item dynamical systems
\item spectral analysis of attention matrices
\end{itemize}

\end{frame}


\begin{frame}{Advanced transformer interpretation (14)}

Transformers can be analyzed through:

\begin{itemize}
\item kernel methods
\item graph neural network analogies
\item dynamical systems
\item spectral analysis of attention matrices
\end{itemize}

\end{frame}


\begin{frame}{Advanced transformer interpretation (15)}

Transformers can be analyzed through:

\begin{itemize}
\item kernel methods
\item graph neural network analogies
\item dynamical systems
\item spectral analysis of attention matrices
\end{itemize}

\end{frame}


\begin{frame}{Advanced transformer interpretation (16)}

Transformers can be analyzed through:

\begin{itemize}
\item kernel methods
\item graph neural network analogies
\item dynamical systems
\item spectral analysis of attention matrices
\end{itemize}

\end{frame}


\begin{frame}{Advanced transformer interpretation (17)}

Transformers can be analyzed through:

\begin{itemize}
\item kernel methods
\item graph neural network analogies
\item dynamical systems
\item spectral analysis of attention matrices
\end{itemize}

\end{frame}


\begin{frame}{Advanced transformer interpretation (18)}

Transformers can be analyzed through:

\begin{itemize}
\item kernel methods
\item graph neural network analogies
\item dynamical systems
\item spectral analysis of attention matrices
\end{itemize}

\end{frame}


\begin{frame}{Advanced transformer interpretation (19)}

Transformers can be analyzed through:

\begin{itemize}
\item kernel methods
\item graph neural network analogies
\item dynamical systems
\item spectral analysis of attention matrices
\end{itemize}

\end{frame}


\begin{frame}{Advanced transformer interpretation (20)}

Transformers can be analyzed through:

\begin{itemize}
\item kernel methods
\item graph neural network analogies
\item dynamical systems
\item spectral analysis of attention matrices
\end{itemize}

\end{frame}

\end{document}
